\chapter{Listas: La Estructura de Datos Universal}
\label{chap:lists}

{\huge$\bigstar$} Principiante

En Lisp, las listas son la estructura de datos fundamental. Toda la sintaxis
del lenguaje se reduce a listas (S-expressions), lo que hace que el código
sea datos y los datos sean código.

\section{Celdas Cons}

La celda cons es el bloque constructor de las listas:

\begin{lispcode}
(cons 1 2)  ;; → (1 . 2)  — dotted pair
\end{lispcode}

Una lista propia es una cadena de celdas cons donde la última apunta a nil:

\begin{lispcode}
(cons 1 (cons 2 (cons 3 nil)))  ;; → (1 2 3)
'(1 2 3 4 5)                     ;; notación literal
(list 1 2 3)                     ;; constructor
\end{lispcode}

\section{Operaciones de Acceso}

\begin{lispcode}
(car '(1 2 3))    ;; → 1  (first)
(cdr '(1 2 3))    ;; → (2 3)  (rest)
(cadr '(1 2 3))   ;; → 2  (second)
(caddr '(1 2 3))  ;; → 3  (third)
(nth 2 '(a b c))  ;; → C
\end{lispcode}

\section{Operaciones de Construcción}

\begin{lispcode}
(push 'x lista)          ;; añade x al inicio
(pop lista)              ;; extrae y retorna el primer elemento
(pushnew 'x lista)       ;; añade si no existe
(append '(1 2) '(3 4))   ;; → (1 2 3 4)
(list* 1 2 '(3 4))       ;; → (1 2 3 4)
(make-list 5 :initial-element 0)  ;; → (0 0 0 0 0)
\end{lispcode}

\section{Búsqueda en Listas}

\begin{lispcode}
(member 3 '(1 2 3 4) :test #'=)       ;; → (3 4)
(find 'c '(a b c d) :test #'eq)       ;; → C
(position 3 '(1 2 3 4 5))              ;; → 2
\end{lispcode}

\section{Listas de Asociación (Alists)}

\begin{lispcode}
(defvar *productos*
  '((:manzana . 0.50)
    (:pan . 2.00)
    (:leche . 1.20)))

(assoc :pan *productos*)   ;; → (:PAN . 2.00)
(rassoc 1.20 *productos*)  ;; → (:LECHE . 1.20)
\end{lispcode}

\section{Operaciones de Conjuntos}

\begin{lispcode}
(adjoin 3 '(1 2 3 4))                ;; → (1 2 3 4) (sin duplicados)
(union '(1 2 3) '(2 3 4))            ;; → (1 2 3 4)
(intersection '(1 2 3) '(2 3 4))     ;; → (2 3)
(set-difference '(1 2 3) '(2 3 4))   ;; → (1)
\end{lispcode}

\section{Ejemplos}

\begin{examplebox}[title=Construcción incremental con push/pop]
\begin{lispcode}
(defvar *pila* nil)
(push 10 *pila*)     ;; → (10)
(push 20 *pila*)     ;; → (20 10)
(pop *pila*)         ;; → 20, *pila* → (10)
(pop *pila*)         ;; → 10, *pila* → nil
\end{lispcode}
\end{examplebox}

\begin{examplebox}[title=Base de datos de productos con alists]
\begin{lispcode}
(defvar *db*
  '((:producto-1 :nombre "Laptop" :precio 1200)
    (:producto-2 :nombre "Mouse"  :precio 25)))

(defun buscar-precio (id)
  (let ((entry (assoc id *db*)))
    (when entry
      (getf (cdr entry) :precio))))

(buscar-precio :producto-1)  ;; → 1200
\end{lispcode}
\end{examplebox}

\begin{examplebox}[title=Aplanamiento recursivo]
\begin{lispcode}
(defun flatten (lst)
  (if (null lst)
      nil
      (if (listp (car lst))
          (append (flatten (car lst))
                  (flatten (cdr lst)))
          (cons (car lst) (flatten (cdr lst))))))

(flatten '((1 2) (3 (4 5)) 6))  ;; → (1 2 3 4 5 6)
\end{lispcode}
\end{examplebox}

\begin{examplebox}[title=Zip de dos listas]
\begin{lispcode}
(defun zip (a b)
  (if (or (null a) (null b))
      nil
      (cons (cons (car a) (car b))
            (zip (cdr a) (cdr b)))))

(zip '(1 2 3) '(a b c))  ;; → ((1 . A) (2 . B) (3 . C))
\end{lispcode}
\end{examplebox}

